% TEMPLATE-MILC-v3.tex - plantilla maestra UNICA de las guias MILC v3.
%
% La usan las tres familias de guias, cada una con su driver:
%   · Grados 6-11  -> scripts/build-guias-g11.py
%   · Bebras       -> scripts/build-guias-bebras.py
%   · Semillero    -> scripts/build-guias-semillero.py
%
% Antes habia tres copias casi identicas (~400 lineas iguales, 6 distintas):
% cada mejora habia que aplicarla tres veces, y en la practica no se hacia
% (el motor de recursos vivia solo en dos de ellas). Lo que varia entre
% familias esta parametrizado en 6 placeholders que cada driver llena
% (se nombran aqui SIN su sintaxis real: el builder reemplaza por texto plano,
% tambien dentro de los comentarios, y un valor multilinea rompe el preambulo):
%
%   HEADER_IZQ        encabezado izquierdo de pagina
%   HEADER_DER        encabezado derecho de pagina
%   PORTADA_ETIQUETA  rotulo sobre el numero grande (GRADO/MOMENTO/MODULO)
%   PORTADA_SERIE     linea de serie de la portada
%   PORTADA_ID        identificador de la guia en la portada
%   PORTADA_CIRCULO   el circulito de grado (°), o vacio si no aplica
%
% El resto de claves las documenta content/guias/_SCHEMA.md. El script valida
% que no queden placeholders sin reemplazar antes de escribir el archivo final.

\documentclass[11pt,letterpaper]{article}
\usepackage[margin=2.0cm]{geometry}
\usepackage[table]{xcolor}
\usepackage{fontspec}
\usepackage[spanish]{babel}
\usepackage{tikz}
% Librerías TikZ para los diagramas de las guías (bloque `recursos:`):
%  · babel  → babel en español vuelve activos caracteres como «>», y TikZ los usa
%             en sus opciones (>=stealth, ->). Sin esta librería, cualquier
%             diagrama con flechas revienta con «\language@active@arg> has an
%             extra }». Debe ir ANTES de que se dibuje nada.
%  · positioning        → sintaxis «below left=2pt and 3pt».
%  · decorations.pathreplacing → llaves ({brace}) para señalar rangos.
\usetikzlibrary{babel,positioning,decorations.pathreplacing}
\usepackage{graphicx}
% Circuitos: el trabajo de electrónica se hace hoy con micro:bit y sus sensores
% integrados (MakeCode, sin cableado), así que no se necesitan esquemáticos.
% Si algún día entran montajes con protoboard, basta descomentar esta línea:
% los diagramas .tex/.tikz declarados en `recursos:` ya se incrustan solos.
% \usepackage{circuitikz}
\usepackage[most]{tcolorbox}
\usepackage{tabularx}
\usepackage{array}
\usepackage{fancyhdr}
\usepackage{titlesec}
\usepackage[document]{ragged2e}  % bandera a la derecha en todo el cuerpo
\usepackage{enumitem}
\usepackage{microtype}

% Helvetica Neue no trae la flecha U+2192, asi que cada «→» escrito literalmente
% en el YAML se compilaba a NADA, en silencio (462 flechas en 61 guias: rutas del
% tipo «paleta → tipografia → lockup» salian rotas). Mapeamos el caracter a su
% version matematica, que si tiene glifo. Asi el autor puede seguir escribiendo
% «→» comodamente en el YAML.
\usepackage{newunicodechar}
\newunicodechar{→}{\ensuremath{\rightarrow}}
% Mismo problema con los simbolos de marca y vinetas: el «✓» de «una palomita ✓»
% salia como cuadrito vacio en el PDF de todas las guias que lo usaban.
\usepackage{pifont}
\newunicodechar{✓}{\ding{51}}
\newunicodechar{✗}{\ding{55}}
\newunicodechar{✔}{\ding{52}}
\newunicodechar{•}{\textbullet}
\newunicodechar{≥}{\ensuremath{\geq}}
\newunicodechar{≤}{\ensuremath{\leq}}

\setmainfont{Helvetica Neue}
\setsansfont{Helvetica Neue}
\setlength{\parindent}{0pt}
\setlength{\parskip}{6pt}
% Sin particion de palabras: en 6.º-7.º una palabra cortada por guion al final
% de linea («Comuni-cacion») frena la lectura mucho mas que una linea corta.
\hyphenpenalty=10000
\exhyphenpenalty=10000
\pretolerance=2000
\tolerance=3000
\setlength{\emergencystretch}{3em}
\linespread{1.10}
\renewcommand{\arraystretch}{1.25}
\setlist{leftmargin=5mm,itemsep=1mm,topsep=1mm}
\newcolumntype{Y}{>{\RaggedRight\arraybackslash}X}

\definecolor{milcMagenta}{HTML}{E6007E}
\definecolor{milcVino}{HTML}{5A0038}
\definecolor{milcTurquesa}{HTML}{0093A5}
\definecolor{milcVerde}{HTML}{58A923}
\definecolor{milcMostaza}{HTML}{E5B400}
\definecolor{milcOcre}{HTML}{B86600}
\definecolor{milcNegro}{HTML}{241816}
\definecolor{milcGris}{HTML}{F7F7F4}
\definecolor{milcLinea}{HTML}{E8E2DD}
\definecolor{milcGradePrimary}{HTML}{84CC16}
\definecolor{milcGradeDark}{HTML}{4D7C0F}
\definecolor{milcGradeSoft}{HTML}{ECFCCB}
\definecolor{milcGradeAccent}{HTML}{CFE7FF}
\definecolor{milcGradeText}{HTML}{FFFFFF}
\color{milcNegro}

\pagestyle{fancy}
\fancyhf{}
\lhead{\small Tecnología e Informática · Bebras}
\rhead{\small Momento 0 · MILC v3}
\cfoot{\small\thepage}
\renewcommand{\headrulewidth}{0pt}

\newtcolorbox{titlebox}[1]{
  enhanced,colback=#1,colframe=#1,arc=7mm,boxrule=0pt,
  left=7mm,right=7mm,top=3mm,bottom=3mm,
  before skip=8pt,after skip=8pt,
  fontupper=\sffamily\bfseries\Large\color{white}
}

\newtcolorbox{softbox}[3]{
  enhanced,colback=#2,colframe=#1,borderline west={4pt}{0pt}{#1},
  arc=2mm,boxrule=.4pt,left=4mm,right=4mm,top=3mm,bottom=3mm,
  before skip=6pt,after skip=8pt,title={#3},coltitle=white,
  colbacktitle=#1,fonttitle=\sffamily\bfseries,fontupper=\RaggedRight,
  attach boxed title to top left={xshift=3mm,yshift=-2mm},
  boxed title style={arc=2mm, boxrule=0pt}
}

\newtcolorbox{drawbox}[2]{
  enhanced,colback=white,colframe=#1,arc=4mm,boxrule=1pt,
  left=5mm,right=5mm,top=4mm,bottom=4mm,before skip=8pt,after skip=8pt,
  title={#2},coltitle=white,colbacktitle=#1,fonttitle=\sffamily\bfseries,
  fontupper=\RaggedRight,
  attach boxed title to top left={xshift=4mm,yshift=-2mm},
  boxed title style={arc=3mm, boxrule=0pt}
}

\newtcolorbox{infoband}[1]{
  enhanced,colback=#1!8,colframe=#1!50,arc=2mm,boxrule=.5pt,
  left=4mm,right=4mm,top=2mm,bottom=2mm,before skip=4pt,after skip=4pt,
  fontupper=\small\RaggedRight
}

% Puente narrativo entre secciones (contrato editorial Conéctate).
% Da continuidad: "Acabas de X. Ahora vas a Y porque Z."
% Visualmente: banda izquierda en color del grado, texto en itálica pequeña.
\newtcolorbox{puentebox}{
  enhanced,colback=milcGradeSoft,colframe=milcGradeDark,
  borderline west={3pt}{0pt}{milcGradeDark},
  arc=0mm,boxrule=0pt,left=5mm,right=4mm,top=2.5mm,bottom=2.5mm,
  before skip=4pt,after skip=8pt,
  fontupper=\itshape\small\color{milcNegro}\RaggedRight
}
% La flecha va en modo matematico a proposito: Helvetica Neue NO tiene el glifo
% U+2192, asi que \textrightarrow se compilaba a NADA («Missing character: There
% is no → in font Helvetica Neue») y el puente salia sin flecha. En modo
% matematico la flecha viene de la fuente matematica y si se dibuja.
\newcommand{\puente}[1]{%
  \begin{puentebox}%
    \textcolor{milcGradeDark}{$\rightarrow$}\hspace{2mm}#1%
  \end{puentebox}%
}

\newcommand{\linead}[1]{\noindent\color{milcLinea}\rule{#1}{0.5pt}\color{milcNegro}}
\newcommand{\respuesta}[1]{\par\vspace{2mm}\foreach \n in {1,...,#1}{\noindent\color{milcLinea}\rule{\linewidth}{0.5pt}\par\vspace{5mm}}\color{milcNegro}}
\newcommand{\checkbox}{\textcolor{milcGradeDark}{\fbox{\phantom{x}}}\hspace{5pt}}

\newcommand{\phasecard}[4]{
\begin{tcolorbox}[enhanced,colback=#3,colframe=#2,arc=3mm,boxrule=.5pt,height=3.05cm,valign=center,halign=center,left=1.5mm,right=1.5mm,top=2mm,bottom=2mm]
{\sffamily\bfseries\fontsize{22}{24}\selectfont\textcolor{#2}{#1}}\par
{\sffamily\bfseries\small #4}
\end{tcolorbox}
}

% ── Piezas del rediseno 2026 (legibilidad 6.º-11.º) ───────────────────
% Banda de estacion: reemplaza el titlebox macizo. Titulo a la izquierda,
% dato practico (tiempo/modalidad) a la derecha, en una sola linea.
\newcommand{\milcestacion}[2]{%
\begin{tcolorbox}[enhanced,colback=#1,colframe=#1,arc=2mm,boxrule=0pt,
  left=5mm,right=5mm,top=2.6mm,bottom=2.6mm,before skip=11pt,after skip=7pt]
{\sffamily\bfseries\fontsize{15}{18}\selectfont\color{white}#2}
\end{tcolorbox}}

% Tarjeta de paso numerada: sustituye las tablas «Pilar / Aplicacion», que
% eran formato de consulta docente, no de estudiante. El numero va en un
% circulo sobre el borde para que la secuencia se lea de un vistazo.
\newcommand{\milcpaso}[3]{%
\begin{tcolorbox}[enhanced,colback=white,colframe=milcLinea,arc=1.5mm,boxrule=.9pt,
  left=14mm,right=4mm,top=3mm,bottom=3mm,before skip=4pt,after skip=4pt,
  fontupper=\RaggedRight,
  overlay={\node[anchor=north west,circle,fill=#2,text=white,inner sep=0pt,
    minimum size=7.4mm,font=\sffamily\bfseries\small]
    at ([xshift=3.6mm,yshift=-3.2mm]frame.north west) {#1};}]
#3
\end{tcolorbox}}

% Casilla marcable de tamano util con lapiz (la anterior era \fbox{\phantom{x}},
% de alto variable segun el texto que la seguia).
\newcommand{\milcmarca}{\framebox[3.8mm]{\rule{0pt}{3.4mm}}}

% Fila marcable de lista suelta (checks de la estacion 1).
\newcommand{\milccheckrow}[1]{%
\par\vspace{1.8mm}\noindent
\makebox[6.4mm][l]{\raisebox{-.55mm}{\textcolor{milcNegro}{\milcmarca}}}%
\parbox[t]{\dimexpr\linewidth-6.4mm\relax}{\RaggedRight #1}\par}

% Celda marcable dentro de la rubrica (tabla de autoevaluacion). Misma
% sangria francesa, pero pensada para vivir dentro de una columna X.
\newcommand{\milcopcion}[1]{%
\makebox[5.2mm][l]{\raisebox{-.55mm}{\textcolor{milcNegro}{\milcmarca}}}%
\parbox[t]{\dimexpr\linewidth-5.2mm\relax}{\RaggedRight #1}}

% ── Recursos visuales de la guía ──────────────────────────────────────
% Declarados en el bloque `recursos:` del YAML y emitidos por el builder.
% \guiaFigura   → imagen rasterizada o PDF (foto, carta celeste, montaje).
% \guiaDiagrama → fuente TikZ/circuitikz (.tex) incrustada como vector:
%                 es la vía para circuitos y esquemáticos editables.
% #1 = ruta relativa al .tex · #2 = pie (puede ir vacío)
\newcommand{\guiaFigura}[2]{%
\begin{center}
\includegraphics[width=0.88\linewidth,height=0.38\textheight,keepaspectratio]{#1}\\[1.5mm]
{\footnotesize\itshape\textcolor{milcNegro!75}{#2}}
\end{center}
}

\newcommand{\guiaDiagrama}[2]{%
\begin{center}
\input{#1}\\[1.5mm]
{\footnotesize\itshape\textcolor{milcNegro!75}{#2}}
\end{center}
}

\begin{document}
\thispagestyle{empty}

% ── Portada ───────────────────────────────────────────────────────────
% Antes: un rectangulo de color a pagina completa. Se veia bien en pantalla,
% pero en la fotocopiadora del colegio se comia el toner y salia gris sucio.
% Ahora la banda ocupa el tercio superior y el resto va en blanco.
\begin{tikzpicture}[remember picture,overlay]
  \path[fill=milcGradePrimary]
    (current page.north west) rectangle ([yshift=-10.6cm]current page.north east);

  \node[anchor=north west,text=milcGradeText] at ([xshift=2.0cm,yshift=-1.55cm]current page.north west)
    {\sffamily\bfseries\fontsize{12}{14}\selectfont MOMENTO};
  \node[anchor=north west,text=milcGradeText,opacity=.85] at ([xshift=2.0cm,yshift=-2.35cm]current page.north west)
    {\sffamily\bfseries\fontsize{8.5}{10}\selectfont BEBRAS · PENSAR COMO UNA MÁQUINA};
  \node[anchor=north east,text=milcGradeText,opacity=.85,align=right] at ([xshift=-2.0cm,yshift=-1.55cm]current page.north east)
    {\sffamily\bfseries\fontsize{9}{12}\selectfont I.E. Sor María Juliana\\Cartago, Valle del Cauca};

  \node[anchor=west,text=milcGradeText] (gradonum) at ([xshift=1.85cm,yshift=-7.1cm]current page.north west)
    {\sffamily\bfseries\fontsize{112}{112}\selectfont 0};
  % El circulito de grado (°) solo aplica a las guias de grado («11°»); en
  % Bebras y Semillero el numero grande es un momento/modulo, no un grado.
  % Se posiciona relativo al nodo `gradonum`, no en coordenadas absolutas.
  

  \node[anchor=west,text=milcGradeText,text width=11.4cm,align=left] at ([xshift=1.5cm,yshift=0.5cm]gradonum.east)
    {
     {\sffamily\bfseries\fontsize{9}{11}\selectfont\MakeUppercase{Curso «Pensar como una máquina» · Pensamiento computacional}}\\[3mm]
     {\sffamily\bfseries\fontsize{21}{25}\selectfont\setlength{\baselineskip}{27pt}El castor que piensa --- las 4 piezas\\[4pt]del pensamiento computacional}\\[3.5mm]
     {\sffamily\bfseries\fontsize{15}{18}\selectfont Bebras · Momento 0}
    };
\end{tikzpicture}

\vspace*{9.4cm}
\vfill

{\sffamily\bfseries\fontsize{8}{10}\selectfont\textcolor{milcGradeDark}{LO QUE TE LLEVAS HOY}}

\begin{tcolorbox}[enhanced,colback=white,colframe=milcNegro,arc=2mm,boxrule=1.6pt,
  left=5mm,right=5mm,top=4mm,bottom=4mm,before skip=3pt,after skip=12pt,
  fontupper=\RaggedRight\fontsize{11}{15.5}\selectfont]
Tu primer reto estilo Bebras resuelto y la tabla «Las 4 piezas en mi día», con la pieza que más te costó reconocer.
\end{tcolorbox}

\begin{tcolorbox}[enhanced,colback=milcGris,colframe=milcGris,arc=2mm,boxrule=0pt,
  left=5mm,right=5mm,top=3.5mm,bottom=3.5mm,after skip=12pt]
{\sffamily\bfseries\fontsize{8}{10}\selectfont\textcolor{milcNegro!55}{ESTA GUÍA ES DE}}\\[3mm]
\begin{tabularx}{\linewidth}{X p{3.2cm} p{3.2cm}}
\linead{\linewidth} & \linead{3.0cm} & \linead{3.0cm}\\[-1mm]
{\fontsize{8.5}{10}\selectfont\textcolor{milcNegro!55}{Nombre completo}} &
{\fontsize{8.5}{10}\selectfont\textcolor{milcNegro!55}{Grupo}} &
{\fontsize{8.5}{10}\selectfont\textcolor{milcNegro!55}{Fecha}}\\
\end{tabularx}
\end{tcolorbox}

\vfill

\noindent\textcolor{milcLinea}{\rule{\linewidth}{1pt}}\\[2mm]
{\sffamily\bfseries\fontsize{9.5}{12}\selectfont Autor: Dr. Álvaro Cárdenas Orozco}\\[1mm]
{\sffamily\fontsize{8.5}{11}\selectfont\textcolor{milcNegro!55}{Metodología MILC · apertura ancestral · pensamiento computacional · triángulo Dussel–estoicismo–Floridi}}

\newpage

\section*{Apertura: El castor que piensa: las 4 piezas del pensamiento computacional}

\begin{softbox}{milcGradeDark}{milcGradeSoft}{Saber ancestral}
En las veredas del Valle del Cauca, cuando una tarea es demasiado grande para una sola familia --- levantar un techo, limpiar una acequia tapada, recoger la cosecha antes de la lluvia ---, se convoca una \textbf{minga}. La comunidad llega un sábado al amanecer y, sin que nadie reparta un manual, el trabajo se organiza solo: unos cargan, otros amarran, otros cocinan el sancocho para todos. El grupo \textbf{parte la tarea enorme en partes pequeñas}, repite lo que funcionó la vez pasada y sigue un orden que todos entienden. Al caer la tarde, lo que era imposible para uno ya está hecho entre muchos. La minga no necesita computadores: es la prueba viva de que resolver algo grande es, sobre todo, \textbf{saber organizar partes pequeñas}.
\end{softbox}

\begin{softbox}{milcTurquesa}{milcGris}{Saber contemporáneo}
\textbf{Bebras} --- «castor» en lituano --- es un reto internacional de \textbf{pensamiento computacional} que cada año reúne a millones de estudiantes en más de 50 países: 45 minutos, unas 15 tareas cortas que se resuelven \textbf{con la cabeza, no con código}. El pensamiento computacional es la forma de resolver problemas que usa un computador --- y que tú ya usabas en la minga ---. Tiene \textbf{cuatro piezas}: descomponer (partir lo grande en partes), reconocer patrones (ver lo que se repite), abstraer (quedarse con lo esencial) y crear algoritmos (escribir pasos exactos). El castor levanta represas que resisten crecientes sin estudiar ingeniería: observa, prueba y organiza. Lo importante no es la respuesta, sino \emph{cómo} la pensaste.
\end{softbox}

\begin{softbox}{milcMostaza}{milcGris}{Pregunta-puente}
\textit{En la minga, la comunidad resuelve un problema enorme partiéndolo en partes y siguiendo un orden acordado. ¿Y si te dijera que una máquina «piensa» exactamente igual --- y que tú también lo haces, sin darte cuenta?}
\end{softbox}

\begin{softbox}{milcVerde}{milcGris}{Saber y saber hacer}
Al terminar podrás: (1) \textbf{identificar} las 4 piezas del pensamiento computacional --- descomposición, patrones, abstracción y algoritmos ---; (2) \textbf{explicar} con tus palabras qué hace cada pieza, con un ejemplo de tu vida; (3) \textbf{analizar} una tarea cotidiana para reconocer en ella las cuatro piezas; (4) \textbf{aplicar} las cuatro piezas para resolver tu primer reto estilo Bebras.
\end{softbox}



\small
\begin{tabularx}{\textwidth}{>{\bfseries}p{3.0cm}Y}
\rowcolor{milcGradePrimary!12}Autor & Dr. Álvaro Cárdenas Orozco\\
DBA & Desarrolla el pensamiento computacional --- descomposición, patrones, abstracción y algoritmos --- para resolver problemas (transversal 6°--11°, alineado a los lineamientos MEN de Tecnología e Informática)\\
Referentes & Valentina Dagienė (Bebras) · Pensamiento computacional (Wing) · Dussel · Estoicismo · Floridi · Saberes del Valle y el Pacífico\\
Producto final & Tu primer reto estilo Bebras resuelto y la tabla «Las 4 piezas en mi día», con la pieza que más te costó reconocer.\\
\end{tabularx}
\normalsize

\puente{Ya sabes que la minga organiza lo grande en partes. Ahora mira el mapa de la sesión: vas a recorrer esas mismas cuatro piezas, una por una, hasta resolver tu primer reto.}

\milcestacion{milcGradeDark}{Ruta MILC: así vamos a aprender}
\begin{tabularx}{\textwidth}{>{\centering\arraybackslash}X>{\centering\arraybackslash}X>{\centering\arraybackslash}X>{\centering\arraybackslash}X}
\phasecard{1}{milcVerde}{milcGris}{Escucha\\[-1mm]\small Observo, escucho y pregunto.} &
\phasecard{2}{milcTurquesa}{milcGris}{Entender\\[-1mm]\small Sistematizo y ordeno.} &
\phasecard{3}{milcMagenta}{milcGris}{Hacer\\[-1mm]\small Construyo y aplico.} &
\phasecard{4}{milcVino}{milcGris}{Cierre\\[-1mm]\small Evalúo y reflexiono.}
\end{tabularx}


\puente{Antes de nombrar las piezas, conviene verlas en acción. Empieza por observar una minga o una tarea de tu casa: las cuatro ya están ahí, escondidas.}

\milcestacion{milcVerde}{Estación 1 de 3 · Escucha}

\begin{softbox}{milcVerde}{milcGris}{Escena para escuchar}
\textbf{Observa una minga} --- o recuerda una en la que hayas estado: un grupo arreglando una casa, organizando un bazar, preparando un sancocho para muchos. Fíjate \emph{cómo} se organizan sin un jefe que dé cada orden. ¿Quién hace qué? ¿Qué repiten porque «así funcionó la vez pasada»? ¿Qué detalles ignoran para no perderse? ¿En qué orden hacen las cosas? Sin saberlo, el grupo ya usa las cuatro piezas del pensamiento computacional.
\end{softbox}

{\sffamily\bfseries\fontsize{8}{10}\selectfont\textcolor{milcNegro!55}{MARCA CUANDO LO HAYAS HECHO}}
\milccheckrow{Nombraste al menos dos tareas pequeñas en que se partió la tarea grande (descomposición).}
\milccheckrow{Señalaste algo que el grupo repitió porque «así funcionó antes» (patrón).}
\milccheckrow{Identificaste un orden de pasos que el grupo siguió (algoritmo).}

\begin{infoband}{milcVerde}


\textbf{En tu cuaderno (Actividad 1 · IDENTIFICA).} Título: «Actividad 1 --- La minga que pienso». \textbf{Formato:} lista de 4 a 6 viñetas. \textbf{Extensión:} una viñeta por cosa observada (quién hace qué, qué se repite, qué se ignora, en qué orden). \textbf{Sabes que terminaste cuando} tu lista nombra, en lenguaje cotidiano, al menos tres de las cuatro piezas en la minga.
\end{infoband}

\puente{Acabas de ver las piezas en lo cotidiano. Ahora vamos a nombrarlas y entenderlas una a una, porque con ellas vas a resolver todo lo que sigue.}

\milcestacion{milcTurquesa}{Estación 2 de 3 · Entender}

\begin{softbox}{milcTurquesa}{milcGris}{Lo que vas a entender}
El pensamiento computacional es la manera de resolver problemas que usa un computador --- y que tú ya usabas en la minga ---. No es magia ni requiere código: son \textbf{cuatro piezas} que, juntas, convierten un problema imposible en algo que sí se puede hacer. Aquí las nombramos una a una, con ejemplos de tu barrio.
\end{softbox}

\milcpaso{1}{milcTurquesa}{\textbf{Partir lo grande en partes pequeñas} que sí puedes resolver. Levantar el techo se vuelve «cargar las tejas», «amarrar las vigas», «subir la escalera». En Bebras, una tarea que parece imposible se vuelve fácil al partirla.}
\milcpaso{2}{milcTurquesa}{\textbf{Ver lo que se repite} para no resolver lo mismo dos veces. Si el bus de la ruta 2 siempre pasa a los minutos 00 y 30, ya sabes cuándo salir sin mirar el horario.}
\milcpaso{3}{milcTurquesa}{\textbf{Quedarte con lo esencial} e ignorar el resto. Cuando dibujas cómo llegar a tu casa, no dibujas cada ladrillo: solo las calles y las esquinas que importan.}
\milcpaso{4}{milcTurquesa}{\textbf{Escribir los pasos exactos}, en el orden correcto. La receta del sancocho es un algoritmo: si cambias el orden o saltas un paso, el resultado cambia.}

\begin{drawbox}{milcTurquesa}{Cómo se piensa una tarea Bebras}
\textbf{1. Lee sin afán} y subraya solo lo que importa (abstracción). \\[1mm] \textbf{2. Parte} el problema en preguntas pequeñas (descomposición). \\[1mm] \textbf{3. Busca} algo que se repita (patrones). \\[1mm] \textbf{4. Arma} los pasos hasta la respuesta (algoritmo). \\[1mm] \textbf{5. Si te trabas, pasa} a otra y vuelve después. En Bebras no gana quien calcula más rápido, sino quien \emph{piensa con más claridad}.
\end{drawbox}

\begin{softbox}{milcMostaza}{milcGris}{Errores comunes y su corrección}
\textbf{Querer resolver todo de un golpe:} el problema parece imposible solo porque no lo has partido. \textbf{Confundir abstraer con «resumir mal»:} abstraer es \emph{elegir} qué dato importa, no quitar al azar. \textbf{Creer que hay que programar:} Bebras se piensa, no se codifica.
\end{softbox}

\begin{infoband}{milcTurquesa}


\textbf{En tu cuaderno (Actividad 2 · EXPLICA).} Título: «Actividad 2 --- Las 4 piezas con mis palabras». \textbf{Formato:} tabla de 2 columnas (pieza · ejemplo de mi día). \textbf{Extensión:} las cuatro piezas, cada una con un ejemplo distinto a los del texto. \textbf{Sabes que terminaste cuando} un compañero entendería tus ejemplos sin que se los expliques.
\end{infoband}

\puente{Ya distingues las cuatro piezas. Es hora de usarlas: vas a resolver tu primer reto estilo Bebras aplicándolas a propósito.}

\milcestacion{milcMagenta}{Estación 3 de 3 · Hacer}

\begin{softbox}{milcMagenta}{milcGris}{Cómo vas a construir tu producto}
Las piezas no son adorno: se usan. Vas a resolver tu \textbf{primer reto estilo Bebras} aplicando las cuatro a propósito. Lee con calma --- la clave está en pensar por partes, no en adivinar.
\end{softbox}

\milcpaso{1}{milcMagenta}{\textbf{Descomposición:} parte el enunciado en datos y pregunta. ¿Qué te dan? ¿Qué te piden?}
\milcpaso{2}{milcMagenta}{\textbf{Patrones:} busca algo que se repita --- una secuencia, un ciclo, una regla.}
\milcpaso{3}{milcMagenta}{\textbf{Abstracción:} ignora los datos de adorno; quédate con los que cambian la respuesta.}
\milcpaso{4}{milcMagenta}{\textbf{Algoritmo:} escribe, en orden, los pasos que te llevan del enunciado a la respuesta.}

\begin{drawbox}{milcMagenta}{El reto: el robot repartidor}
\checkbox \textbf{Reto (Nivel A).} Un robot mira al \textbf{norte}. Recibe estas órdenes, en orden: \emph{avanza, gira a la derecha, avanza, gira a la derecha, avanza}. ¿Hacia dónde mira al final? \\[2mm]
\checkbox Marca: \quad \fbox{\phantom{x}}~Norte \quad \fbox{\phantom{x}}~Sur \quad \fbox{\phantom{x}}~Este \quad \fbox{\phantom{x}}~Oeste \\[2mm]
\checkbox Escribe \textbf{qué pieza} usaste para resolverlo y por qué.
\end{drawbox}

\begin{softbox}{milcMagenta}{milcGris}{Plantilla de trabajo}
\textbf{Resuélvelo paso a paso (algoritmo).} Empiezas mirando al \textbf{norte}. Cada «gira a la derecha» cambia el rumbo 90°: norte $\to$ este $\to$ sur. Los «avanza» mueven al robot, pero \emph{no} cambian el rumbo. Por eso, al final, el robot mira al \textbf{sur}. Pieza usada: \textbf{algoritmo} --- seguir pasos exactos, en orden.
\end{softbox}

\begin{infoband}{milcMagenta}


\textbf{En tu cuaderno (Actividad 3 · APLICA).} Título: «Actividad 3 --- Mi primer reto Bebras». \textbf{Formato:} respuesta marcada + 3 renglones de explicación. \textbf{Extensión:} la respuesta y por qué, nombrando la pieza usada. \textbf{Sabes que terminaste cuando} tu explicación convencería a alguien que aún no sabe la respuesta.
\end{infoband}

\puente{Resolviste el reto. Ahora deja tu evidencia: el reto explicado y la tabla de las piezas en tu día, para mostrar que sabes pensar así.}

\milcestacion{milcMostaza}{Producto de la sesión}
\textbf{Tu primer reto Bebras --- resuelto y explicado --- más la tabla «Las 4 piezas en mi día»}.
\begin{softbox}{milcMostaza}{milcGris}{Criterios de calidad}
(1) Resolviste el reto y marcaste la respuesta correcta (el sur). (2) Explicaste con tus palabras cómo lo pensaste, nombrando al menos una de las cuatro piezas. (3) Tu tabla «Las 4 piezas en mi día» tiene un ejemplo propio por cada pieza. (4) Distingues claramente descomponer de abstraer, sin confundirlos. (5) Tu trabajo muestra que entendiste que Bebras se piensa, no se programa.
\end{softbox}

\puente{Terminaste el producto. Detente a mirar qué cambió en ti --- no la nota, sino tu manera de pensar.}

\milcestacion{milcVino}{Evaluación liberadora · cinco dimensiones}

\begin{softbox}{milcVino}{milcGris}{Sentido de la evaluación}
Evaluar no es solo revisar una nota. Es reconocer qué aprendí, qué cambió en mí, qué emociones manejé, cómo me relaciono con la ciudad y la comunidad, y qué le dejaré a quien viene detrás.
\end{softbox}

\textbf{Mi reflexión liberadora en cinco dimensiones.} Completa cada frase iniciada:

\small
\begin{tabularx}{\textwidth}{>{\bfseries\RaggedRight\arraybackslash}p{3.4cm}Y>{\RaggedRight\arraybackslash}p{4.6cm}}
\rowcolor{milcVino}\color{white}Dimensión & \color{white}\bfseries Mi respuesta (frame starter) & \color{white}\bfseries Algo que descubrí\\
1. Personal & Lo que más cambió en mí fue \linead{6.5cm} & \linead{4.2cm}\\
2. Emocional & La emoción más fuerte fue \linead{2.5cm} y la regulé \linead{3cm} & \linead{4.2cm}\\
3. Ciudadana & En democracia esto importa porque \linead{6cm} & \linead{4.2cm}\\
4. Local (Cartago) & En mi barrio sería útil para \linead{6.5cm} & \linead{4.2cm}\\
5. Intergeneracional & A mi abuela/abuelo le diría \linead{6.5cm} & \linead{4.2cm}\\
\end{tabularx}
\normalsize

\begin{infoband}{milcVino}
\textbf{Para la próxima sustentación.} Vuelve a esta tabla en seis meses. Verás que las respuestas cambian — no porque lo aprendido cambie, sino porque tú cambias. La evaluación liberadora es una conversación contigo a través del tiempo.
\end{infoband}


\puente{Cerramos pensando en grande: tres voces para entender qué significa, para ti y para tu comunidad, aprender a pensar como una máquina.}

\milcestacion{milcGradeDark}{Triángulo de pensamiento · cierre filosófico}

\begin{softbox}{milcVino}{milcGris}{Dussel · lente del nosotros}
\textit{El saber del pueblo también es ciencia.}

La minga ya organizaba el trabajo como un algoritmo --- partir, repetir, ordenar --- mucho antes de que existiera la palabra «computador». El pensamiento computacional no llegó de afuera: tu comunidad ya lo practicaba.

\textbf{¿Qué saberes de mi familia o mi comunidad ya eran «pensar como una máquina», aunque nadie los llamara así?}
\end{softbox}
\linead{\linewidth}\\[1mm]
\linead{\linewidth}

\begin{softbox}{milcMostaza}{milcGris}{Estoicismo · Marco Aurelio · cuidado interior}
\textit{Tienes poder sobre tu mente, no sobre los hechos externos. Date cuenta de esto y hallarás la fuerza.}

Cuando un problema se ve enorme y el tiempo corre, no controlas el reloj, pero sí tu manera de enfrentarlo: respira, pártelo en pasos pequeños y empieza por el primero.

\textbf{¿Puedo calmarme y partir un problema difícil en pasos pequeños, en vez de bloquearme?}
\end{softbox}
\linead{\linewidth}\\[1mm]
\linead{\linewidth}

\begin{softbox}{milcTurquesa}{milcGris}{Floridi · lente de la infoesfera}
\textit{Somos inforgs: organismos informacionales que habitan la infoesfera.}

Pensar es procesar información paso a paso. Aprender las cuatro piezas es entender cómo piensan las máquinas con las que ya convives --- y cómo piensas tú.

\textbf{¿En qué se parece mi mente a la de un computador, y en qué es distinta?}
\end{softbox}
\linead{\linewidth}\\[1mm]
\linead{\linewidth}

\begin{infoband}{milcNegro}
\textbf{Nota para el docente.} Las tres formulaciones del triángulo son la síntesis del autor de la guía, no citas textuales de los pensadores. Si vas a citarlos en clase, remite a la obra original.
\end{infoband}

\milcestacion{milcVino}{Mi autoevaluación · marca una casilla por fila}

{\small
\setlength{\extrarowheight}{2.2pt}
\begin{tabularx}{\textwidth}{>{\RaggedRight\arraybackslash\bfseries}p{3.0cm}YYY}
\rowcolor{milcVino}\color{white}\bfseries Criterio & \color{white}\bfseries Lo logré & \color{white}\bfseries En proceso & \color{white}\bfseries Necesito apoyo\\[.6mm]
Comprensión del tema & \milcopcion{Explico las ideas con mis palabras.} & \milcopcion{Reconozco las ideas, falta precisión.} & \milcopcion{Necesito repasar lo básico.}\\[.8mm]
\rowcolor{milcGris}Pensamiento computacional & \milcopcion{Usé los pilares al armar mi producto.} & \milcopcion{Usé algunos pilares.} & \milcopcion{Necesito releer la estación 2.}\\[.8mm]
Aplicación y producto & \milcopcion{Mi producto cumple los criterios.} & \milcopcion{Está armado, con ajustes pendientes.} & \milcopcion{Aún está en construcción.}\\[.8mm]
\rowcolor{milcGris}Reflexión 5 dimensiones & \milcopcion{Respondí cada una con honestidad.} & \milcopcion{Respondí algunas con detalle.} & \milcopcion{Mis respuestas son muy generales.}\\[.8mm]
Triángulo de pensamiento & \milcopcion{Conecté las 3 voces con mi trabajo.} & \milcopcion{Conecté al menos una.} & \milcopcion{No las relaciono con mi proceso.}\\[.8mm]
\end{tabularx}}

\vspace{3mm}
\textbf{Hoy entendí que:}\\[1mm]
\linead{\linewidth}\\[1mm]
\linead{\linewidth}

\textbf{Mi compromiso para la próxima sesión es:}\\[1mm]
\linead{\linewidth}\\[1mm]
\linead{\linewidth}

\begin{softbox}{milcMostaza}{milcGris}{Cita de despedida · Marco Aurelio}
\textit{``El obstáculo en el camino se vuelve el camino. Lo que estorba al camino se vuelve camino.''}\\[1mm]
Cada error de hoy, bien observado, se convierte en pieza del camino que pisarás mañana.
\end{softbox}

\small
\begin{tabularx}{\textwidth}{>{\bfseries}p{3.6cm}Y}
\rowcolor{milcGradeDark!12}Nombre completo: & \linead{10cm}\\
Fecha: & \linead{4cm} \quad Curso/grupo: \linead{4cm}\\
Mi firma: & \linead{9cm}\\
\end{tabularx}
\normalsize

\begin{infoband}{milcGradeDark}
\textbf{Para la próxima sesión.} Caza las cuatro piezas en una tarea de tu casa --- cocinar, alistar el morral, organizar tu cuarto: pártela en pasos, marca dónde usas un patrón y dónde abstraes, y escríbelo. Trae un ejemplo para compartir con el grupo.
\end{infoband}

\end{document}
